\chapter{Joint Replacement}

\section*{Introduction \& Clinical Indications}
Joint replacement, medically termed arthroplasty, is a definitive surgical procedure involving the meticulous excision of diseased or damaged articular cartilage and underlying subchondral bone from a severely affected joint. This is followed by the precise implantation of prosthetic components, typically crafted from advanced biocompatible materials such as metal alloys, high-density polyethylene, and ceramics. These implants are engineered to meticulously replicate the natural joint's intricate anatomy and biomechanics, aiming to restore smooth, pain-free motion and stability. This transformative intervention is primarily undertaken to alleviate intractable chronic pain, correct debilitating deformities, and significantly improve functional mobility in patients suffering from end-stage degenerative (e.g., osteoarthritis) or inflammatory (e.g., rheumatoid arthritis) joint diseases, most commonly affecting the hip, knee, and shoulder. The overarching goal is to provide a durable, functional articulation that profoundly enhances the patient's quality of life and independence.

\begin{itemize}
  \item Arthroplasty
  \item Total Joint Arthroplasty (TJA)
  \item Total Joint Reconstruction
  \item Prosthetic Joint Replacement
  \item Endoprosthetic Replacement
  \item Hemiarthroplasty (for partial joint replacement, e.g., hip)
\end{itemize}

The fundamental principle of joint replacement surgery is to meticulously resurface the articulating ends of a severely damaged joint, thereby eliminating the primary source of pain and restoring optimal biomechanical function. This is achieved by precisely resecting the diseased or damaged articular cartilage and a small amount of underlying subchondral bone, which are the main contributors to friction, inflammation, and chronic pain. The resected bone ends are then meticulously prepared and shaped using specialized instruments to create a precise fit for the engineered prosthetic components.

These prostheses typically comprise a metallic femoral or humeral component, a metallic or polyethylene tibial or glenoid component, and a highly durable polyethylene bearing surface that articulates between them. The components are secured to the bone either through biological fixation (press-fit, allowing for bone ingrowth into porous surfaces) or with polymethylmethacrylate (PMMA) bone cement, a decision often guided by patient age, bone quality, and surgeon preference. The new artificial joint is designed to provide a smooth, low-friction articulation, facilitating a wide range of motion without pain. The technical goals of the procedure include achieving proper limb alignment, restoring joint stability, balancing soft tissue tension around the joint, and ensuring optimal implant positioning to maximize longevity and functional outcome, ultimately aiming to improve the patient's quality of life significantly.

The concept of replacing diseased joints dates back to the late 19th century, with early attempts at arthroplasty focusing on interpositional techniques using various biological materials like fascia, muscle, or even pig bladder to create a new joint surface. Themistocles Gluck, a German surgeon, is widely credited with performing some of the earliest documented joint replacements in the 1890s, utilizing ivory components fixed with nickel-plated screws and glue. These pioneering efforts, while innovative for their time, were largely unsuccessful due to high rates of infection and limitations of the materials used.

Significant advancements began to emerge in the mid-20th century. In the 1940s, the Judet brothers introduced acrylic femoral head prostheses, which offered some short-term relief but often failed due to aseptic loosening. The true revolution in modern joint replacement, however, began in the 1960s with the groundbreaking work of Sir John Charnley in the United Kingdom. Charnley pioneered the low-friction total hip replacement, combining a small-diameter metallic femoral head articulating with a high-density polyethylene acetabular cup, meticulously fixed with bone cement. His meticulous surgical technique, emphasis on ultra-sterile operating environments (including filtered air systems), and profound understanding of biomechanical principles laid the robust foundation for the widespread success of contemporary total hip arthroplasty. This success subsequently influenced the rapid development and acceptance of other joint replacements, such as total knee replacement, which gained widespread clinical adoption in the 1970s, and shoulder arthroplasty in later decades, continually evolving with material science and surgical techniques.

Joint replacement surgery is indicated for patients experiencing severe, debilitating joint pathology that significantly impairs their quality of life and has proven refractory to comprehensive conservative management.

\begin{itemize}
  \item Severe, chronic, and debilitating joint pain that is unresponsive to non-surgical interventions, profoundly impacting daily activities and sleep.
  \item Advanced osteoarthritis (degenerative joint disease) characterized by extensive cartilage loss, bone-on-bone friction, osteophyte formation, and structural damage visible on imaging.
  \item Rheumatoid arthritis and other inflammatory arthropathies leading to progressive joint destruction, severe pain, and significant joint deformity.
  \item Avascular necrosis (osteonecrosis) of the femoral head, humeral head, or other joint surfaces, resulting in bone collapse, severe pain, and functional impairment.
  \item Post-traumatic arthritis following severe intra-articular fractures or complex joint injuries that have led to irreparable joint damage and secondary osteoarthritis.
  \item Significant joint stiffness and a profound loss of active and passive range of motion that severely compromises functional independence and mobility.
  \item Uncorrectable joint deformity, such as severe varus or valgus angulation of the knee, causing gait abnormalities, instability, and progressive pain.
  \item Failed previous joint preservation surgeries, osteotomies, or less extensive arthroplasty procedures that have not provided adequate long-term relief or function.
\end{itemize}

Joint replacement is predominantly considered a primary, definitive surgical treatment for end-stage joint disease. It is typically positioned as the ultimate therapeutic option when a comprehensive array of non-surgical interventions has been exhausted and failed to provide adequate pain relief or functional improvement. These conservative measures include physical therapy, activity modification, weight loss, oral analgesics, non-steroidal anti-inflammatory drugs (NSAIDs), and intra-articular injections (e.g., corticosteroids, hyaluronic acid).

In certain acute scenarios, such as displaced femoral neck fractures in elderly patients, total hip arthroplasty or hemiarthroplasty may be considered a primary, immediate treatment to facilitate early mobilization, reduce pain, and mitigate the high risks of complications associated with prolonged immobility. However, for the vast majority of degenerative conditions, it serves as the final and most effective option after a structured trial of conservative management. In cases of failed prior arthroplasty due to aseptic loosening, periprosthetic infection, or implant wear, revision joint replacement surgery is performed, which would then be classified as a secondary or tertiary procedure, often with increased complexity and potential risks compared to a primary replacement.

\section*{Relevant Surgical Anatomy}
Joint replacement surgery necessitates a profound understanding of the specific joint's anatomy, including its bony architecture, capsuloligamentous structures, surrounding musculature, and critical neurovascular bundles. For a typical synovial joint, the articular surfaces are covered by hyaline cartilage, which provides a low-friction gliding surface. Beneath this lies the subchondral bone, which provides structural support. The joint is encased by a fibrous capsule, reinforced by intrinsic and extrinsic ligaments that provide stability. Surrounding muscles facilitate movement and dynamic stability.

In the hip, the primary components are the femoral head and the acetabulum, a cup-shaped socket formed by the ilium, ischium, and pubis. Key ligaments include the iliofemoral, pubofemoral, and ischiofemoral ligaments. The major muscles involved in surgical approaches include the gluteal muscles (maximus, medius, minimus), piriformis, and iliopsoas. Critical neurovascular structures at risk include the femoral artery and vein, the sciatic nerve (posterior approach), and the femoral nerve (anterior approach). The medial and lateral circumflex femoral arteries are vital for femoral head vascularity. For the knee, the distal femoral condyles articulate with the tibial plateau, with the patella gliding in the trochlear groove. Stability is provided by the anterior and posterior cruciate ligaments (ACL, PCL) and the medial and lateral collateral ligaments (MCL, LCL). The quadriceps and hamstring muscles are crucial for function. The popliteal artery and vein, and the common peroneal nerve (lateral aspect of the knee), are particularly vulnerable during surgical dissection. Lymphatic drainage generally follows major vascular bundles, with regional nodes in the groin for the hip and popliteal fossa for the knee. Surgical landmarks, such as the anterior superior iliac spine for hip approaches or the patella and tibial tubercle for knee approaches, guide incision placement and component alignment.

\section*{Preoperative Workup \& Preparation}
Thorough preoperative workup and preparation are paramount to optimize patient outcomes, minimize surgical risks, and ensure a smooth recovery. This comprehensive process typically involves:

\begin{itemize}
  \item \textbf{Medical Evaluation:} A comprehensive medical history and physical examination are performed by the orthopedic surgeon and often by an internist or hospitalist. This includes a review of all comorbidities, medications, and allergies.
  \item \textbf{Laboratory Tests:} Routine blood tests are ordered, including a complete blood count (CBC), electrolyte panel, renal and liver function tests, coagulation profile (PT/INR, PTT), and urinalysis to screen for infection.
  \item \textbf{Cardiac \& Anesthesia Clearance:} Patients, especially those with pre-existing cardiac or pulmonary conditions, often require clearance from a cardiologist and/or an anesthesiologist to assess their fitness for major surgery and anesthesia. An electrocardiogram (ECG) and sometimes a chest X-ray are standard.
  \item \textbf{Imaging Studies:} Preoperative weight-bearing X-rays of the affected joint are essential for surgical planning, templating implant sizes, and assessing bone quality and deformity. Other imaging modalities like MRI or CT scans may be ordered in specific, complex cases.
  \item \textbf{Dental Screening:} A complete dental screening and clearance are highly recommended to rule out any occult dental infections, which could potentially seed the new joint prosthesis and lead to a devastating periprosthetic joint infection.
  \item \textbf{Medication Adjustment:} Patients are typically instructed to discontinue certain medications, such as blood thinners (anticoagulants and antiplatelet agents), several days to a week prior to surgery to minimize the risk of intraoperative and postoperative bleeding. Other medications may require dosage adjustments.
  \item \textbf{NPO Status:} Patients must adhere strictly to NPO (nil per os -- nothing by mouth) guidelines for a specified period (typically 6-8 hours) before surgery to prevent aspiration during anesthesia induction.
  \item \textbf{Preoperative Education \& Physical Therapy (Prehab):} Patients receive detailed education about the surgical procedure, expected recovery trajectory, potential risks, and the importance of rehabilitation. Preoperative physical therapy, or "prehab," can help strengthen muscles, improve range of motion, and educate patients on postoperative exercises, potentially aiding a faster recovery.
  \item \textbf{Antibiotic Prophylaxis:} Prophylactic broad-spectrum antibiotics are administered intravenously shortly before the surgical incision to significantly reduce the risk of surgical site infection.
  \item \textbf{Informed Consent:} The patient must provide fully informed consent after a thorough discussion of the procedure, its anticipated benefits, potential risks, alternative treatments, and the expected recovery process.
\end{itemize}

Joint replacement surgery, while highly effective, has specific contraindications and requires careful precautions to ensure patient safety and optimal outcomes.

\textbf{Situations requiring optimization, infection assessment, or an alternative plan:}
\begin{itemize}
  \item Active systemic infection (e.g., sepsis, bacteremia) or active local joint infection (e.g., septic arthritis, osteomyelitis), as this dramatically increases the risk of periprosthetic joint infection, which is a catastrophic complication.
  \item Severe neuropathic joint disease (Charcot joint) with progressive bone destruction and instability, which can lead to early implant loosening and failure due to lack of protective sensation and abnormal biomechanics.
  \item Inadequate bone stock or severe, uncorrectable osteoporosis that prevents stable and durable fixation of the prosthetic components.
  \item Uncontrolled, severe medical comorbidities (e.g., decompensated heart failure, recent myocardial infarction, uncontrolled diabetes mellitus, severe pulmonary hypertension) that pose an unacceptably high anesthetic or surgical risk.
  \item Paralysis or severely inadequate muscle function around the joint that would preclude functional recovery and lead to persistent instability or inability to rehabilitate.
  \item Skin breakdown or severe dermatological conditions directly over the planned surgical incision site, increasing infection risk.
\end{itemize}

\textbf{Relative Contraindications \& Precautions:}
\begin{itemize}
  \item Morbid obesity (Body Mass Index > 40 kg/m\textasciicircum{}2), which increases surgical complexity, anesthetic risks, and the risk of complications such as infection, deep vein thrombosis (DVT), pulmonary embolism (PE), and implant loosening.
  \item Significant peripheral vascular disease or severe skin conditions (e.g., psoriasis, venous stasis ulcers) around the surgical site, which can impair wound healing and increase infection risk.
  \item History of prior joint infection, requiring careful evaluation, extensive preoperative workup, and potentially a two-stage revision procedure.
  \item Active urinary tract infection or other distant infections that should be aggressively treated and cleared prior to elective surgery.
  \item Unrealistic patient expectations regarding surgical outcomes, or a demonstrated history of non-compliance with medical advice or postoperative rehabilitation protocols.
  \item Young age, where the long-term longevity of the implant may be a concern, and alternative joint-preserving procedures might be considered first, weighing the risk of future revision surgeries.
  \item Severe psychiatric illness or cognitive impairment that would prevent the patient from understanding the procedure or participating in rehabilitation.
\end{itemize}

\section*{Surgical Technique \& Procedure}
Joint replacement surgery is a highly specialized procedure requiring meticulous planning and execution. The general steps are as follows:

\begin{itemize}
  \item \textbf{Anesthesia:} The procedure is typically performed under general anesthesia, often supplemented with regional anesthesia (e.g., spinal or epidural block, peripheral nerve block) to provide superior postoperative pain control and reduce the need for systemic opioids.
  \item \textbf{Patient Positioning:} The patient is carefully positioned on the operating table to optimize surgical access to the affected joint while ensuring patient safety and comfort. Positioning varies by joint and surgical approach (e.g., supine or lateral decubitus for hip, supine for knee, beach chair or lateral for shoulder).
  \item \textbf{Surgical Exposure \& Incision:} After sterile preparation and draping of the surgical field, an incision is made through the skin and subcutaneous tissues. The surgeon then meticulously dissects through muscle layers and opens the joint capsule to expose the diseased articular surfaces. The length and location of the incision depend on the specific joint and the chosen surgical approach (e.g., anterior, anterolateral, posterior for hip; medial parapatellar for knee; deltopectoral for shoulder).
  \item \textbf{Bone Resection \& Preparation:} Specialized surgical instruments, including oscillating saws, reamers, and broaches, are used to precisely resect the damaged bone and cartilage from the articulating surfaces (e.g., femoral head and acetabulum for hip; distal femur and proximal tibia for knee; humeral head and glenoid for shoulder). The remaining bone ends are then meticulously shaped and reamed to create a precise, stable bed for the prosthetic components, often guided by intraoperative jigs and alignment guides to ensure accurate limb alignment and component orientation.
  \item \textbf{Trial Reduction \& Soft Tissue Balancing:} Trial prosthetic components are temporarily inserted to assess joint stability, range of motion, and limb length equality (especially for hip replacement). The surgeon carefully balances the surrounding soft tissues (ligaments, capsule) to ensure optimal tension, preventing instability or excessive tightness.
  \item \textbf{Implant Implantation:} The definitive prosthetic components are then implanted. This may involve cementing the components into place using polymethylmethacrylate bone cement, or press-fitting them into the bone to allow for biological ingrowth into porous surfaces. For total hip replacement, this includes the acetabular cup, femoral stem, and femoral head. For total knee replacement, it involves the femoral component, tibial component, and often a patellar component.
  \item \textbf{Final Assessment \& Irrigation:} Once the definitive components are securely in place, the joint is reduced, and the surgeon performs a final assessment of stability, range of motion, and alignment. The surgical site is thoroughly irrigated with sterile saline to remove any bone debris or blood clots.
  \item \textbf{Closure:} The joint capsule, muscle layers, and subcutaneous tissues are meticulously closed with absorbable sutures. The skin is closed with sutures, staples, or surgical glue. A surgical drain may be placed temporarily to collect any excess fluid or blood, typically removed within 24-48 hours.
\end{itemize}

\section*{Complications \& Risks}
While joint replacement surgery is generally safe and highly effective, like any major surgical procedure, it carries potential risks and complications that patients must be aware of.

\begin{itemize}
  \item \textbf{Anesthesia-related risks:}
    \begin{itemize}
      \item Adverse reactions to anesthetic agents, including allergic reactions.
      \item Respiratory complications such as pneumonia, atelectasis, or exacerbation of pre-existing lung disease.
      \item Cardiovascular events, including myocardial infarction, stroke, or arrhythmias.
      \item Nerve damage from positioning or regional blocks, usually temporary.
    \end{itemize}
  \item \textbf{General surgical risks:}
    \begin{itemize}
      \item \textbf{Bleeding:} Intraoperative or postoperative hemorrhage, potentially requiring blood transfusion, which carries its own risks.
      \item \textbf{Infection:} Surgical site infection, ranging from superficial wound infection to deep periprosthetic joint infection (PJI), a serious complication that may necessitate prolonged antibiotic treatment, surgical debridement, or even implant removal and revision surgery.
      \item \textbf{Blood clots:} Deep vein thrombosis (DVT) in the legs, which can potentially dislodge and travel to the lungs, causing a life-threatening pulmonary embolism (PE). Prophylactic anticoagulation is routinely used.
      \item \textbf{Wound healing problems:} Delayed wound healing, dehiscence, or hematoma formation.
    \end{itemize}
  \item \textbf{Procedure-specific risks:}
    \begin{itemize}
      \item \textbf{Joint dislocation:} Particularly in total hip replacement, the ball of the new joint can come out of the socket, often requiring reduction and sometimes revision surgery.
      \item \textbf{Implant loosening or wear:} Over time, the prosthetic components can loosen from the bone (aseptic loosening) or the bearing surfaces can wear out, necessitating revision surgery.
      \item \textbf{Nerve or blood vessel injury:} Damage to adjacent nerves (e.g., sciatic nerve in hip replacement, common peroneal nerve in knee replacement, axillary nerve in shoulder replacement) can cause weakness, numbness, or paralysis. Damage to blood vessels can lead to bleeding, pseudoaneurysm, or ischemia.
      \item \textbf{Periprosthetic fracture:} A fracture occurring around the implant, either during surgery (intraoperative) or postoperatively due to trauma or stress, often requiring additional surgical fixation or revision.
      \item \textbf{Leg length discrepancy:} More common in hip replacement, where one leg may feel longer or shorter than the other, potentially requiring shoe lifts or further intervention.
      \item \textbf{Persistent pain or stiffness:} Despite technically successful surgery, some patients may experience ongoing pain, limited range of motion, or stiffness, sometimes requiring manipulation under anesthesia or further surgery.
      \item \textbf{Heterotopic ossification:} Abnormal bone formation in the soft tissues around the joint, which can restrict motion and cause pain.
      \item \textbf{Allergic reaction to implant materials:} Although rare, some patients may develop sensitivity or allergic reactions to metal components (e.g., nickel, cobalt, chromium).
      \item \textbf{Patellar complications (for knee replacement):} Patellar fracture, patellar maltracking, or avascular necrosis of the patella.
    \end{itemize}
\end{itemize}

\section*{Postoperative Care \& Recovery}
Immediately after joint replacement surgery, the patient is transferred to the Post-Anesthesia Care Unit (PACU) for close monitoring of vital signs, pain levels, and the surgical site. Pain management is a critical priority, often involving a multimodal approach with oral and intravenous analgesics, nerve blocks, or patient-controlled analgesia (PCA). Early mobilization is a cornerstone of modern joint replacement recovery, with many patients encouraged to sit up, stand, and even take a few steps with assistance (e.g., walker, crutches) on the day of surgery or the following day, depending on the joint and patient condition.

Over the first 24-48 hours, the patient typically remains in the hospital. Physical therapy commences almost immediately, focusing on gentle range-of-motion exercises, muscle strengthening, and gait training with assistive devices. Anticoagulant medications are prescribed to prevent deep vein thrombosis. The surgical drain, if placed, is usually removed within 24-48 hours. The length of hospital stay varies but is often 1-3 days for total knee or hip replacement, with some patients undergoing outpatient surgery. During the first 1-2 weeks post-discharge, patients focus on wound care, managing pain with oral medications, and adhering to a structured home exercise program or outpatient physical therapy. Activity restrictions are in place (e.g., weight-bearing precautions, hip precautions to prevent dislocation). Diet is advanced as tolerated. Long-term recovery continues for several months, with continued improvement in strength, flexibility, and function over the first year. Most patients can resume light recreational activities by 3-6 months and achieve maximal functional recovery by 12-18 months.

During joint replacement surgery, the surgeon meticulously documents the macroscopic findings within the joint. These typically include severe erosion or complete absence of articular cartilage, exposing eburnated (polished, sclerotic) subchondral bone. Other common findings are osteophyte formation (bone spurs) at the joint margins, subchondral cysts, and areas of bone collapse, particularly in cases of avascular necrosis. The synovium may appear inflamed, hypertrophied, or fibrotic, reflecting chronic degenerative or inflammatory processes. The integrity of ligaments and the joint capsule is also assessed.

Any resected bone and cartilage, as well as synovial tissue, are routinely sent for pathological examination. The pathology report typically confirms the clinical diagnosis, describing features consistent with osteoarthritis (e.g., cartilage fibrillation, chondrocyte necrosis, subchondral bone sclerosis, osteophyte formation, mild synovitis) or inflammatory arthropathy (e.g., synovial hyperplasia, pannus formation, inflammatory cell infiltrates). In cases of avascular necrosis, the report will describe necrotic bone and marrow. If revision surgery is performed, the pathologist may identify evidence of implant wear debris (e.g., polyethylene particles, metallic debris) within the synovium, leading to a foreign body giant cell reaction and aseptic loosening. The absence of acute inflammatory cells or microorganisms is crucial to rule out periprosthetic joint infection.

A normal and successful postoperative course following joint replacement is characterized by a predictable progression of healing and functional recovery. Clinically, the patient experiences progressive pain relief, with acute surgical discomfort gradually diminishing over days to weeks as the incision heals and inflammation subsides. There is a steady and measurable improvement in joint mobility and strength, allowing for increasing weight-bearing and functional activities as guided by the physical therapy regimen. The surgical incision heals cleanly, without signs of infection such as excessive redness, warmth, swelling, purulent discharge, or fever.

On follow-up radiographic imaging (X-rays), the prosthetic components are well-aligned, securely fixed to the bone, and show no evidence of migration, loosening, or periprosthetic fracture. The patient is able to gradually reduce reliance on assistive devices, eventually walking independently. Most patients can resume the majority of their normal daily activities, including light recreational sports and driving, within a few months post-surgery. The ultimate goal is a significant enhancement in their overall quality of life, independence, and ability to participate in desired activities with minimal or no pain, typically reaching maximal improvement within 6-12 months.

Comprehensive postoperative care and structured follow-up are essential for optimizing recovery and ensuring the long-term success of joint replacement surgery.

\begin{itemize}
  \item \textbf{Postoperative Visits:} Regular orthopedic follow-up appointments are scheduled, typically at 2 weeks, 6 weeks, 3 months, 6 months, 1 year, and then annually or every few years thereafter, depending on the surgeon's protocol and patient's progress.
  \item \textbf{Suture/Staple Removal:} Skin sutures or staples are usually removed at the first postoperative visit, approximately 10-14 days after surgery, once the wound has adequately healed.
  \item \textbf{Imaging:} Follow-up X-rays of the replaced joint are performed at key intervals (e.g., 6 weeks, 1 year, 5 years, and then every 5-10 years) to monitor implant position, stability, and detect any early or late signs of loosening, wear, or other complications.
  \item \textbf{Physical Therapy/Rehabilitation:} Continued outpatient physical therapy is crucial for several weeks to months to regain strength, flexibility, and optimize functional recovery. Adherence to a prescribed home exercise program is also vital for long-term success.
  \item \textbf{Activity Resumption Guidance:} The surgeon provides specific instructions on activity restrictions, weight-bearing status, and when it is safe to gradually resume various activities, including driving, work, and recreational sports, tailored to the individual's progress and the specific joint replaced.
  \item \textbf{Warning Signs:} Patients are educated on critical warning signs of potential complications, such as increasing pain, redness, warmth, swelling, fever, chills, purulent drainage from the incision, or a sudden inability to bear weight, and are instructed to contact their surgeon immediately if any of these occur.
\end{itemize}
Long-term surveillance is essential to monitor implant longevity, detect potential late complications like aseptic loosening or infection, and ensure sustained functional outcomes.

\section*{Outcomes \& Prognosis}
Joint replacement surgery, particularly total hip arthroplasty (THA) and total knee arthroplasty (TKA), stands as one of the most common and consistently successful elective surgical procedures performed globally. In the United States, over 1 million THAs and TKAs are performed annually, with projections indicating a substantial increase in volume over the next decade, driven by an aging population, rising rates of obesity, and expanding indications. The incidence of joint replacement generally increases with age, with the majority of procedures performed in individuals over 60 years old, although younger patients with post-traumatic arthritis, inflammatory conditions, or avascular necrosis also undergo surgery. There is a slight female predominance for TKA, while THA incidence is more evenly distributed between sexes. Osteoarthritis is overwhelmingly the most common indication, accounting for over 90\% of all joint replacements. The mortality rate associated with elective joint replacement is remarkably low, typically less than 0.5\%, while major morbidity rates (e.g., infection, DVT, PE, dislocation) range from 1-5\%, varying by patient comorbidities, the specific joint replaced, and surgical approach.

The outcomes and prognosis following joint replacement surgery are generally excellent, with consistently high rates of patient satisfaction and significant improvements in quality of life. For both total hip and total knee arthroplasty, success rates in terms of profound pain relief and functional improvement exceed 90-95\%. Implant survival rates are also very favorable, with 10-year survival rates typically above 90-95\% and 20-year survival rates often exceeding 80-85\% for modern prostheses, as evidenced by large national joint registries. Functional outcomes include significantly improved walking distance, reduced reliance on assistive devices, and the ability to participate in many low-impact activities such as swimming, cycling, and golf.

Prognostic factors influencing long-term success include patient age, bone quality, adherence to postoperative rehabilitation protocols, and the absence of complications such as periprosthetic infection or early aseptic loosening. While recurrence of the original joint disease is not possible as the joint is replaced, complications like periprosthetic infection, aseptic loosening, or implant wear can necessitate revision surgery, which carries higher risks and potentially less favorable outcomes than primary surgery. Overall, joint replacement offers a highly effective and durable solution for end-stage joint disease, allowing the vast majority of patients to return to an active and pain-free lifestyle for many years, profoundly enhancing their independence and well-being.

\section*{References}
\begin{enumerate}
  \item MedlinePlus. Joint Replacement. U.S. National Library of Medicine; 2024. Available from: https://medlineplus.gov/jointreplacement.html
  \item Azar FM, Beaty JH, Canale ST. Campbell's Operative Orthopaedics. 14th ed. Philadelphia, PA: Elsevier; 2021.
  \item American Academy of Orthopaedic Surgeons (AAOS). Total Joint Replacement. 2023. Available from: https://orthoinfo.aaos.org/en/treatment/total-joint-replacement/
  \item National Institute of Arthritis and Musculoskeletal and Skin Diseases (NIAMS). Joint Replacement Surgery. National Institutes of Health; 2023. Available from: https://www.niams.nih.gov/health-topics/joint-replacement-surgery
\end{enumerate}

\clearpage
